\documentclass[10pt,a4paper]{article}
\usepackage[T1]{fontenc}
\usepackage[utf8]{inputenc}

\usepackage[margin=1in,left=0.75in,right=0.7in,top=0.5in,bottom=0.45in]{geometry}
\usepackage{xcolor}
\usepackage{tikz}
\usepackage{parskip}
\usepackage{enumitem}
\usepackage{needspace}
\usepackage{eso-pic}
\usepackage{listings}
\usepackage{array}
\usepackage{tabularx}
\usepackage{booktabs}
\usepackage[colorlinks=true,urlcolor=teal,linkcolor=teal]{hyperref}

\definecolor{ink}{HTML}{111827}
\definecolor{teal}{HTML}{0F766E}
\definecolor{sand}{HTML}{FBFAF6}
\definecolor{rule}{HTML}{D9D6CC}
\definecolor{muted}{HTML}{6B7280}
\definecolor{warn}{HTML}{B45309}
\pagecolor{sand}\color{ink}
\pagestyle{empty}

\AddToShipoutPictureBG{\begin{tikzpicture}[remember picture,overlay]
  \fill[teal] (current page.north west) rectangle ([xshift=0.35in]current page.south west);
\end{tikzpicture}}

\newcommand{\kicker}[1]{{\sffamily\bfseries\small\color{teal}\MakeUppercase{#1}}}
\newcommand{\lbl}[1]{{\sffamily\bfseries\footnotesize\color{muted}\MakeUppercase{#1}}}
\newcommand{\sect}[1]{\needspace{5\baselineskip}\vspace{0.5em}\par\kicker{#1}\par\vspace{-0.5em}{\color{rule}\rule{\linewidth}{0.5pt}}\par\vspace{-0.2em}}
\newcommand{\todo}[1]{{\color{warn}\bfseries[TODO: #1]}}
\setlength{\parskip}{0.3em}
\setlist[itemize]{itemsep=1pt,topsep=1pt,leftmargin=1.2em}

\newcolumntype{Y}{>{\raggedright\arraybackslash\setlength{\parskip}{0pt}}X}
\renewcommand{\arraystretch}{1.25}
\setlength{\tabcolsep}{5pt}

\lstset{
  basicstyle=\ttfamily\footnotesize\color{ink},
  breaklines=true,
  breakatwhitespace=true,
  frame=single,
  rulecolor=\color{rule},
  backgroundcolor=\color{white},
  framesep=5pt,
  xleftmargin=3pt,
  xrightmargin=3pt,
  columns=fullflexible,
  keepspaces=true,
  aboveskip=4pt,
  belowskip=4pt
}

\newcommand{\notes}[4]{%
\lbl{Notes}\par
\begin{itemize}
  \item \textbf{What changed in the prompt:} #1
  \item \textbf{What improved in the output:} #2
  \item \textbf{What still failed:} #3
  \item \textbf{What I would try next:} #4
\end{itemize}}

\begin{document}

\begin{center}
  {\sffamily\bfseries\LARGE Make It Do Something}\\[0.2em]
  {\sffamily\large PF-04 portfolio feature submission}\\[0.4em]
  Azizul Abedin Azmi
\end{center}

\sect{The one feature}

My portfolio has exactly one dynamic feature: a working contact form. A visitor
can enter a name, email address, and message on the live site. Netlify Forms
receives and stores the submission in the Netlify dashboard, so I do not need
to build or pay for a separate server.

\begin{tabularx}{\linewidth}{@{}lY@{}}
\toprule
\textbf{Feature} & Contact form with Netlify Forms\\
\textbf{Live site} & \url{https://azizulabedinazmi.netlify.app/}\\
\textbf{Hosting tier} & Netlify free tier\\
\textbf{Form name} & \texttt{contact}\\
\textbf{Source} & \url{https://github.com/azizulabedinazmi/FlyRank_PF-04}\\
\bottomrule
\end{tabularx}

\sect{What a backend is}

A backend is the part of an application that runs away from the visitor's
browser. It receives requests, validates or processes data, stores data, and
returns a result. The browser shows the interface, but the backend is what
makes a feature persist or perform work after the page is loaded.

For this project, I did not write a custom backend. Netlify provides the
backend service through Netlify Forms. This is useful because the portfolio
is static, but the contact form still needs somewhere reliable to send and
store messages.

\sect{How the data flows}

\begin{enumerate}[leftmargin=1.5em,itemsep=2pt]
  \item A visitor opens the portfolio over HTTPS and fills in the contact
        form.
  \item The form has the name \texttt{contact} and the \texttt{netlify}
        attribute. The hidden \texttt{form-name} field identifies the form.
  \item The browser sends the form fields to Netlify when the visitor presses
        ``Send message''. Required fields prevent an empty name, invalid email,
        or empty message from being submitted.
  \item Netlify receives the request, matches it to the detected
        \texttt{contact} form, and stores the submission in the Forms
        dashboard.
  \item I can review the message in Netlify and reply using the visitor's
        submitted email address.
\end{enumerate}

\begin{lstlisting}[language=HTML]
<form name="contact" method="POST" netlify
      netlify-honeypot="bot-field">
  <input type="hidden" name="form-name" value="contact">
  <input type="text" name="name" required>
  <input type="email" name="email" required>
  <textarea name="message" required></textarea>
</form>
\end{lstlisting}

\sect{Why this meets the brief}

\begin{itemize}
  \item \textbf{Exactly one feature:} the contact form is the only
        server-backed feature on the portfolio.
  \item \textbf{Free tier:} the site and form are hosted by Netlify without a
        paid server or database.
  \item \textbf{End to end:} the form is deployed at the live URL and Netlify
        Form detection is enabled for the site.
  \item \textbf{Understanding:} the data flow is explained in plain words
        above rather than treating the service as a black box.
\end{itemize}

\sect{Evidence of the real test}

\todo{Insert a screenshot showing the submitted test message in Netlify
Forms, and optionally the confirmation page or notification email.}

Test details to record:

\begin{tabularx}{\linewidth}{@{}lY@{}}
\toprule
\textbf{Test URL} & \url{https://azizulabedinazmi.netlify.app/}\\
\textbf{Test date} & \todo{Add date}\\
\textbf{Test name} & \todo{Add the name used}\\
\textbf{Result} & \todo{Confirm the submission appears in Netlify Forms}\\
\bottomrule
\end{tabularx}

\sect{Reflection}

\notes
{I changed the form from a custom redirect and custom JavaScript submission
flow to Netlify's documented native \texttt{netlify} form pattern.}
{The implementation is smaller and uses the platform's built-in form
detection, submission handling, and dashboard storage.}
{The repository cannot prove that a real visitor submission reached the
dashboard; that requires one final manual test in Netlify.}
{Submit a test message, capture the dashboard entry, and add that screenshot
to this report before uploading it.}

\end{document}
